\clearpage\begingroup
\SourceLeaf{318}\NoSourceHead\HistoricalFoot{XI1.}{40}
\begin{center}
{\large [BESTIMMUNGEN DES SONNEN-\\DURCHMESSERS AM MITTAGSFERNROHR.]}\par\medskip
\textls[150]{BRIEFWECHSEL.}\par\medskip
\rule{.64\textwidth}{.3pt}\par
{\small [Briefwechsel zwischen \textsc{Gauss} und \textsc{Bessel}, Leipzig 1880.]}\par
\rule{.64\textwidth}{.3pt}\par\medskip
\textsc{Gauss} an \textsc{Bessel}, Göttingen, 5. März 1820.
\end{center}

\noindent\hspace*{\parindent}Da man an dem \textsc{Reichenbach}schen Meridiankreise, wenn man nichts übereilen will, jeden Mittag nicht gut mehr als einen Sonnenrand, d. i. den obern oder den untern beobachten kann, so war mir daran gelegen zu sehen, einen wie grossen Sonnendurchmesser ein ähnliches Fernrohr gibt, und welche Uebereinstimmung dabei stattfindet; ich habe daher neulich alle meine mit dem [\textsc{Reichenbach}schen] Mittagsfernrohr beobachteten Sonnendurchmesser genau berechnet und folgendes gefunden:
\begin{center}
\begin{tabular}{@{}rlrrcc@{}}
&&&Fäden&Durchg.-Zeit&Reducirtes Resultat\\
1819&Aug.&29.&3&\(129\overset{\mathrm{s}}{,}17\)&\(128\overset{\mathrm{s}}{,}23\)\\
&Dec.&26.&5&\(142{,}52\)&\(128{,}21\)\\
1820&Jan.&1.&2&\(142{,}30\)&\(128{,}31\)\\
&&4.&7&\(141{,}91\)&\(128{,}22\)\\
&&12.&7&\(140{,}71\)&\(128{,}13\)\\
&&13.&7&\(140{,}46\)&\(128{,}06\)\\
&&14.&7&\(140{,}60\)&\(128{,}35\)\\
&&15.&7&\(140{,}59\)&\(128{,}49\)\\
&&24.&7&\(138{,}39\)&\(128{,}15\)\\
&Febr.&1.&7&\(136{,}44\)&\(128{,}01\)\\
&&9.&7&\(134{,}60\)&\(128{,}02\)\\
&&15.&7&\(133{,}56\)&\(128{,}28\)\\
&&16.&7&\(133{,}16\)&\(128{,}09\)\\
&&20.&7&\(132{,}53\)&\(128{,}27\)\\
&&22.&7&\(131{,}86\)&\(127{,}98\)\\
&&27.&7&\(131{,}17\)&\(128{,}16\)\\
&&28.&7&\(130{,}95\)&\(128{,}11\)\\
&&29.&7&\(130{,}71\)&\(128{,}02\)\\
&März&5.&7&\(130{,}10\)&\(128{,}11\)\\\cline{2-6}
&\multicolumn{4}{l}{Mittel aus 122 Fäden}&\(128{,}160.\)
\end{tabular}
\end{center}
\clearpage\SourceLeaf{319}\HistoricalHead{314}{L}{PRAKT. UND SPHÄR. ASTRONOMIE. REICHENBACHSCHES PASSAGENINSTRUMENT.}
\noindent Dies stimmt sehr nahe mit der Angabe in \textsc{Carlini}s Tafeln. Ausgeschlossen sind hiebei die Beobachtungen im Herbst, weil Anfang September das Sonnenglas einen Sprung bekommen hatte, und ich so lange, bis ich von \textsc{Reichenbach} ein neues erhielt, gezwungen war, mich eines andern zu bedienen, wodurch die Sonne roth und nicht so schön begrenzt erschien, wie durch die \textsc{Reichenbach}sch[en]. Die Beobachtungen sind aus eben diesem Grunde, weil ich ihnen weniger traute, auch nicht zahlreich,
\begin{center}21 Fäden geben im Mittel \(128\overset{\mathrm{s}}{,}213\).\end{center}
\noindent An den Tagen, wo die Differenz vom Mittel besonders gross ist, wie am 14. und 15. Januar, wallte die Sonne sehr stark.

\begin{center}\textsc{Bessel} an \textsc{Gauss}, Königsberg, 20. November 1820.\end{center}
\noindent\hspace*{\parindent}. . . . . . \{Den mittlern Halbmesser der Sonne geben mir die oben angeführten 65 Sonnenbeobachtungen im Sinne der AR \(=16'1\overset{\prime\prime}{,}245\) und im Sinne der Decl[ination \(=\)] \(16'1\overset{\prime\prime}{,}11\); die erste Zahl ist kaum von der verschieden, welche Sie mir einmal mittheilten.\}

\begin{center}\textsc{Gauss} an \textsc{Bessel}, Göttingen, 11. März 1821.\end{center}
\noindent\hspace*{\parindent}. . . . . . Von meinen eignen Beobachtungen am M[ittags-]F[ernrohr] geben 271 Appulse von Januar—Julius 1820 den mittlern Halbmesser [der Sonne] \(16'1\overset{\prime\prime}{,}01\).

\begin{center}BEMERKUNGEN.\end{center}
\begingroup\small
\noindent\hspace*{\parindent}Das \textsc{Reichenbach}sche Passageninstrument erhielt \textsc{Gauss} im November 1818. Die im Nachlass erhaltenen dürftigen und zum Abdruck nicht geeigneten Notizen umfassen Beobachtungen von etwa 60 Hauptsternen, deren Deklinationen \textsc{Harding} gleichzeitig am \textsc{Repsold}schen Meridiankreise bestimmte (vgl. Werke IV, S. 426). Diese Notizen weisen im wesentlichen auf eine Beobachtungsreihe vom 27. Juli bis 29. August 1819 hin, obwohl die Beobachtungen wenigstens bis zum März 1820 fortgesetzt wurden, wie aus dem vorstehenden Briefe hervorgeht. Veröffentlicht sind jedoch nur die gelegentlichen Anschlüsse von Planeten (Werke VI, S. 423—424).

Ausserdem diente das Passageninstrument in den Jahren 1819—1822 zur Beobachtung von Mondsternen, an die nach der bekannten von \textsc{Nicolai} angegebenen Methode auf \textsc{Gauss}’ Vorschlag von ihm selbst, von \textsc{Nicolai} in Mannheim, \textsc{Soldner} in Bogenhausen, \textsc{Encke} auf dem Seeberg, \textsc{Schumacher} in Altona, F. G. W. \textsc{Struve} in Dorpat und \textsc{Bessel} in Königsberg der Mond angeschlossen wurde (vgl. Werke VI, S. 441 und sonst).

Man sehe auch den Aufsatz über \textsc{Gauss}’ astronomische Arbeiten, Werke XI2.
\begin{flushright}\textsc{Brendel}.\end{flushright}
\begin{center}\rule{.24\textwidth}{.3pt}\end{center}
\par\endgroup
\clearpage\endgroup
